\section{HeatGeo}
\label{sec:method}

\subsection{Teacher Semantic Graph}

HeatGeo first encodes the full training corpus with the teacher and constructs a teacher semantic graph. For each example $x_i$, let
\begin{equation}
    \mathcal{N}_k^T(i)
    =
    \operatorname{TopK}_{j\neq i}\cos(t_i,t_j)
\end{equation}
be its top-$k$ teacher nearest neighbors. To reduce noisy one-directional edges and hubness, HeatGeo uses a mutual $k$NN graph:
\begin{equation}
    j \in \mathcal{M}_k^T(i)
    \Longleftrightarrow
    j \in \mathcal{N}_k^T(i)
    \;\mathrm{and}\;
    i \in \mathcal{N}_k^T(j).
\end{equation}
The weighted adjacency matrix $W^T$ is
\begin{equation}
    W_{ij}^{T}
    =
    \begin{cases}
    \exp(\cos(t_i,t_j)/\tau_g), & j \in \mathcal{M}_k^T(i),\\
    0, & \mathrm{otherwise},
    \end{cases}
\end{equation}
where $\tau_g$ controls graph sharpness. Let $D^T$ be the diagonal degree matrix. The row-normalized transition matrix is
\begin{equation}
    P^T = (D^T)^{-1}W^T.
\end{equation}
Each row $P_i^T$ defines a one-step semantic random-walk distribution from anchor $x_i$.

\subsection{Multi-Scale Heat Diffusion Targets}

A one-step transition captures direct neighbors, but embedding spaces also contain higher-order semantic structure. HeatGeo therefore uses multiple diffusion scales:
\begin{equation}
    P_r^T = (P^T)^r,
    \qquad r \in \mathcal{R},
\end{equation}
where $\mathcal{R}$ is typically $\{1,2,4\}$. Small $r$ preserves sharp local neighborhoods, while larger $r$ captures semantic regions and multi-hop communities. In implementation, diffusion is approximated sparsely by propagating probability mass through graph neighbor lists and retaining top candidates.

For each anchor $x_i$, HeatGeo builds a candidate set
\begin{equation}
    C_i = C_i^{+} \cup C_i^{h} \cup C_i^{r},
\end{equation}
where $C_i^{+}$ contains high-probability diffusion neighbors, $C_i^{h}$ contains teacher hard negatives, and $C_i^{r}$ contains random negatives. The teacher diffusion target at scale $r$ is the diffusion mass restricted and renormalized over $C_i$:
\begin{equation}
    p_{i,r}^{T}(j)
    =
    \frac{(P_r^T)_{ij}}
    {\sum_{j' \in C_i}(P_r^T)_{ij'}},
    \qquad j\in C_i.
\end{equation}
This target is a soft distribution over semantic relatedness, rather than a binary positive-negative label.

\subsection{Student Neighborhood Distribution}

Given an anchor embedding $s_i$ and candidate embeddings $\{s_j:j\in C_i\}$, the student induces a distribution over the same candidate set:
\begin{equation}
    p_i^{S}(j)
    =
    \frac{\exp(\cos(s_i,s_j)/\tau_s)}
    {\sum_{j'\in C_i}\exp(\cos(s_i,s_{j'})/\tau_s)},
    \qquad j\in C_i,
\end{equation}
where $\tau_s$ is a student temperature. The main diffusion distillation loss is
\begin{equation}
    \mathcal{L}_{\mathrm{diff}}
    =
    \frac{1}{|B|}
    \sum_{i\in B}
    \sum_{r\in\mathcal{R}_{e}}
    \omega_r
    D_{\mathrm{KL}}
    \left(
        p_{i,r}^{T} \| p_i^{S}
    \right),
\end{equation}
where $B$ is a mini-batch, $\omega_r$ is the scale weight, and $\mathcal{R}_{e}$ is the set of active scales at epoch $e$.

\subsection{Spectral Manifold Anchoring}

Diffusion matching transfers local and multi-hop neighborhoods, but a student can still fit local candidates while distorting broader community structure. HeatGeo therefore computes spectral coordinates from the teacher graph. Let the symmetric normalized graph Laplacian be
\begin{equation}
    L^T = I - (D^T)^{-1/2}W^T(D^T)^{-1/2}.
\end{equation}
Let $U^T\in\mathbb{R}^{N\times m}$ contain the first $m$ non-trivial eigenvectors of $L^T$, and let $u_i^T$ be the row corresponding to $x_i$. A learnable projection maps student embeddings to this spectral space:
\begin{equation}
    z_i^S = W_{\mathrm{spec}}s_i.
\end{equation}
The spectral objective is
\begin{equation}
    \mathcal{L}_{\mathrm{spec}}
    =
    \frac{1}{|B|}
    \sum_{i\in B}
    \left\|
    z_i^S-u_i^T
    \right\|_2^2.
\end{equation}
This term encourages the student to retain global teacher topology without requiring hidden states or teacher forward passes during student training.

\subsection{Training Objective and Curriculum}

HeatGeo includes a lightweight pointwise anchor for stability:
\begin{equation}
    \mathcal{L}_{\mathrm{anchor}}
    =
    \frac{1}{|B|}
    \sum_{i\in B}
    \left(1-\cos(W_a s_i,t_i)\right),
\end{equation}
where $W_a$ maps student embeddings to the teacher dimension. When paired views are available, we also use a contrastive task loss $\mathcal{L}_{\mathrm{task}}$ implemented as InfoNCE over two student encodings of each example pair.

The final objective is
\begin{equation}
\begin{aligned}
    \mathcal{L}_{\mathrm{HeatGeo}}
    ={}&
    \lambda_{\mathrm{diff}}\mathcal{L}_{\mathrm{diff}}
    +
    \lambda_{\mathrm{spec}}\mathcal{L}_{\mathrm{spec}}
    \\
    &+
    \lambda_{\mathrm{anchor}}\mathcal{L}_{\mathrm{anchor}}
    +
    \lambda_{\mathrm{task}}\mathcal{L}_{\mathrm{task}}.
\end{aligned}
\end{equation}
We use a diffusion-scale curriculum: the first epoch uses one-step diffusion, the second epoch adds two-step diffusion, and later epochs add the full scale set and spectral anchoring. This avoids forcing long-range, smoother structure before the student has learned reliable local neighborhoods.

\begin{algorithm}[t]
\caption{HeatGeo Training}
\label{alg:heatgeo}
\begin{algorithmic}[1]
\REQUIRE Corpus $\mathcal{D}$, teacher $T$, student $S_{\theta}$, graph size $k$, diffusion scales $\mathcal{R}$
\STATE Cache teacher embeddings $t_i=T(x_i)$ for all $x_i\in\mathcal{D}$
\STATE Build a mutual $k$NN teacher graph $W^T$
\STATE Compute transition matrix $P^T=(D^T)^{-1}W^T$
\STATE Precompute sparse diffusion candidates and targets $p_{i,r}^T$
\STATE Compute spectral coordinates $u_i^T$ from the graph Laplacian
\FOR{each training epoch}
    \STATE Activate diffusion scales according to the curriculum
    \FOR{each mini-batch $B$}
        \STATE Encode anchors and candidate texts with $S_{\theta}$
        \STATE Compute student candidate distributions $p_i^S$
        \STATE Compute $\mathcal{L}_{\mathrm{diff}}$, $\mathcal{L}_{\mathrm{spec}}$, $\mathcal{L}_{\mathrm{anchor}}$, and $\mathcal{L}_{\mathrm{task}}$
        \STATE Update $\theta$ by minimizing $\mathcal{L}_{\mathrm{HeatGeo}}$
    \ENDFOR
\ENDFOR
\RETURN Distilled student encoder $S_{\theta}$
\end{algorithmic}
\end{algorithm}
